
\documentclass{article}
\usepackage{colortbl}
\usepackage{makecell}
\usepackage{multirow}
\usepackage{supertabular}

\begin{document}

\newcounter{utterance}

\centering \large Interaction Transcript for game `clean\_up', experiment `1\_medium\_7obj\_bg', episode 3 with mistral{-}large{-}3{-}azure.
\vspace{24pt}

{ \footnotesize  \setcounter{utterance}{1}
\setlength{\tabcolsep}{0pt}
\begin{supertabular}{c@{$\;$}|p{.15\linewidth}@{}p{.15\linewidth}p{.15\linewidth}p{.15\linewidth}p{.15\linewidth}p{.15\linewidth}}
    \# & $\;$A & \multicolumn{4}{c}{Game Master} & $\;\:$B\\
    \hline

    \theutterance \stepcounter{utterance}  
    & & \multicolumn{4}{p{0.6\linewidth}}{
        \cellcolor[rgb]{0.9,0.9,0.9}{
            \makecell[{{p{\linewidth}}}]{
                \texttt{\tiny{[P1$\langle$GM]}}
                \texttt{Аз съм вашият ръководител на играта и вие играете съвместна игра със следната мрежа като игрално поле:} \\
\\ 
\texttt{```} \\
\texttt{ 1234567} \\
\texttt{╔═╤═════╗} \\
\texttt{║◌│◌A◌◌◌║ 1} \\
\texttt{╟─┤◌T◌◌◌║ 2} \\
\texttt{║O└─┐◌◌◌║ 3} \\
\texttt{║◌◌◌├─┐◌║ 4} \\
\texttt{║◌◌S│◌│◌║ 5} \\
\texttt{╟─┬─┴─┼─╢ 6} \\
\texttt{║P│◌DM│◌║ 7} \\
\texttt{╚═╧═══╧═╝} \\
\\ 
\texttt{```} \\
\\ 
\texttt{* Горният ръб показва x{-}координати, които нарастват надясно, а десният ръб показва y{-}координати, които нарастват надолу.} \\
\texttt{* Следните обекти са разположени на случаен принцип върху вашата мрежа: 'P', 'O', 'T', 'S', 'D', 'A', 'M'.} \\
\\ 
\\ 
\texttt{Другият играч вижда вариант на игралното поле, където обектите са разположени на различни произволни места. Вие не можете да видите полето на другия играч, а той не може да види вашето.} \\
\\ 
\\ 
\texttt{**Цел:**} \\
\\ 
\texttt{И двамата играчи трябва да преместят обектите на своя фон така, че еднакви обекти да се окажат на едни и същи координати. Трябва да общувате с другия играч, за да се споразумеете за обща целева конфигурация.} \\
\\ 
\\ 
\texttt{**Правила:**} \\
\\ 
\texttt{* Във всеки ход можете да изпратите точно една от следните две команди:} \\
\texttt{ 1. `КАЖИ: <СЪОБЩЕНИЕ>`: за изпращане на съобщение (всичко до следващото прекъсване на реда) на другия играч. Аз ще го препратя на вашия партньор.} \\
\texttt{2. `ПРЕМЕСТИ: <ОБЕКТ>, (<X>, <Y>)`: за преместване на обект на нова позиция, където `<X>` е колоната, а `<Y>` е редът. Ще ви информирам дали преместването ви е успешно или не.} \\
\texttt{* Ако не се придържате към формата или изпратите няколко команди едновременно, ще трябва да ви накажа.} \\
\texttt{* Ако и двамата играчи натрупат повече от 16 наказания, и двамата губите играта.} \\
\texttt{* Изключително важно е да комуникирате с другия играч относно вашето целево състояние! Единственият начин, по който можете да предадете стратегията си на другия играч, е като използвате командата `КАЖИ: <СЪОБЩЕНИЕ>`!} \\
\\ 
\\ 
\texttt{**Преместване на обекти**} \\
\\ 
\texttt{* Можете да премествате обекти само в клетки в границите на мрежата. Целевата клетка трябва да е празна, т.е. да съдържа само символа '◌'.} \\
\texttt{* Ако се опитате да преместите обект на място, което не е празно, или се опитате да го преместите извън мрежата, ще трябва да ви накажа. Получавате още един опит.} \\
\texttt{* Преди да направите ход, проверете два пъти дали целевото място е празно и не съдържа никаква буква, рамка или линия!} \\
\\ 
\\ 
\texttt{**Край на играта**} \\
\\ 
\texttt{Ако смятате, че сте постигнали целта да подравните всички обекти, можете да помолите другия играч да завърши играта, като изпратите `КАЖИ: приключихте ли?`. Ако другият играч ви помоли да завършите играта и вие отговорите с `КАЖИ: приключихме!`, играта ще приключи.} \\
\\ 
\texttt{И двамата играчи печелят, ако играта приключи в рамките на 28 рунда, където един рунд се определя като двама играчи, всеки от които изпраща валидна команда.} \\
\\ 
\texttt{**Оценяване:**} \\
\\ 
\texttt{Колкото по{-}близо са еднаквите обекти в двете игрални полета, толкова повече точки получавате. Наказанията намаляват точките ви. Можете ли да подобрите рекорда?} \\
\\ 
\\ 
\texttt{Моля, изпратете съобщение на другия играч, за да започнете играта!} \\
            }
        }
    }
    & \\ \\

    \theutterance \stepcounter{utterance}  
    & \multicolumn{4}{p{0.6\linewidth}}{
        \cellcolor[rgb]{0.9,0.9,0.9}{
            \makecell[{{p{\linewidth}}}]{
                \texttt{\tiny{[P1$\rangle$GM]}}
                \texttt{Да започнем с координирането! Ето моите текущи позиции на обектите:} \\
\\ 
\texttt{{-} P: (1,7)} \\
\texttt{{-} O: (1,3)} \\
\texttt{{-} T: (3,2)} \\
\texttt{{-} S: (3,5)} \\
\texttt{{-} D: (5,7)} \\
\texttt{{-} A: (3,1)} \\
\texttt{{-} M: (6,7)} \\
\\ 
\texttt{Моля, споделете вашите позиции, за да можем да изработим обща целева конфигурация. Предлагам да започнем с подреждане на обектите в низходящ ред по азбучен ред (A, D, M, O, P, S, T) и да ги разположим в горната част на полето (редове 1{-}3), за да избегнем пресичане с линиите.} \\
\\ 
\texttt{Какво мислите? Можете ли да споделите вашите позиции?} \\
            }
        }
    }
    & & \\ \\

    \theutterance \stepcounter{utterance}  
    & & & \multicolumn{2}{p{0.3\linewidth}}{
        \cellcolor[rgb]{0.9,0.9,0.9}{
            \makecell[{{p{\linewidth}}}]{
                \texttt{\tiny{[GM$|$GM]}}
                \texttt{Да започнем с координирането! Ето моите текущи позиции на обектите:} \\
\\ 
\texttt{{-} P: (1,7)} \\
\texttt{{-} O: (1,3)} \\
\texttt{{-} T: (3,2)} \\
\texttt{{-} S: (3,5)} \\
\texttt{{-} D: (5,7)} \\
\texttt{{-} A: (3,1)} \\
\texttt{{-} M: (6,7)} \\
\\ 
\texttt{Моля, споделете вашите позиции, за да можем да изработим обща целева конфигурация. Предлагам да започнем с подреждане на обектите в низходящ ред по азбучен ред (A, D, M, O, P, S, T) и да ги разположим в горната част на полето (редове 1{-}3), за да избегнем пресичане с линиите.} \\
\\ 
\texttt{Какво мислите? Можете ли да споделите вашите позиции?} \\
            }
        }
    }
    & & \\ \\

    \theutterance \stepcounter{utterance}  
    & & & \multicolumn{2}{p{0.3\linewidth}}{
        \cellcolor[rgb]{0.9,0.9,0.9}{
            \makecell[{{p{\linewidth}}}]{
                \texttt{\tiny{[GM$|$GM]}}
                \texttt{Invalid response: Да започнем с координирането! Ето моите текущи позиции на обектите:} \\
\\ 
\texttt{{-} P: (1,7)} \\
\texttt{{-} O: (1,3)} \\
\texttt{{-} T: (3,2)} \\
\texttt{{-} S: (3,5)} \\
\texttt{{-} D: (5,7)} \\
\texttt{{-} A: (3,1)} \\
\texttt{{-} M: (6,7)} \\
\\ 
\texttt{Моля, споделете вашите позиции, за да можем да изработим обща целева конфигурация. Предлагам да започнем с подреждане на обектите в низходящ ред по азбучен ред (A, D, M, O, P, S, T) и да ги разположим в горната част на полето (редове 1{-}3), за да избегнем пресичане с линиите.} \\
\\ 
\texttt{Какво мислите? Можете ли да споделите вашите позиции?} \\
            }
        }
    }
    & & \\ \\

    \theutterance \stepcounter{utterance}  
    & & \multicolumn{4}{p{0.6\linewidth}}{
        \cellcolor[rgb]{0.9,0.9,0.9}{
            \makecell[{{p{\linewidth}}}]{
                \texttt{\tiny{[P1$\langle$GM]}}
                \texttt{Наказание: Трябва да започнете играта, като изпратите съобщение на другия играч!} \\
\\ 
\texttt{Уверете се, че вашият отговор съдържа само `КАЖИ: <СЪОБЩЕНИЕ>` или `ПРЕМЕСТИ: <ОБЕКТ>, (<X>, <Y>)` и нищо друго!} \\
\\ 
\texttt{Колективно сте натрупали 1 от 16 наказания.} \\
\texttt{Моля, опитайте отново!} \\
            }
        }
    }
    & \\ \\

    \theutterance \stepcounter{utterance}  
    & \multicolumn{4}{p{0.6\linewidth}}{
        \cellcolor[rgb]{0.9,0.9,0.9}{
            \makecell[{{p{\linewidth}}}]{
                \texttt{\tiny{[P1$\rangle$GM]}}
                \texttt{КАЖИ: Здравей! Ето моите обекти и техните позиции: P(1,7), O(1,3), T(3,2), S(3,5), D(5,7), A(3,1), M(6,7). Какви са вашите? Да изберем обща цел!} \\
            }
        }
    }
    & & \\ \\

    \theutterance \stepcounter{utterance}  
    & & & \multicolumn{2}{p{0.3\linewidth}}{
        \cellcolor[rgb]{0.9,0.9,0.9}{
            \makecell[{{p{\linewidth}}}]{
                \texttt{\tiny{[GM$|$GM]}}
                \texttt{КАЖИ: Здравей! Ето моите обекти и техните позиции: P(1,7), O(1,3), T(3,2), S(3,5), D(5,7), A(3,1), M(6,7). Какви са вашите? Да изберем обща цел!} \\
            }
        }
    }
    & & \\ \\

    \theutterance \stepcounter{utterance}  
    & & \multicolumn{4}{p{0.6\linewidth}}{
        \cellcolor[rgb]{0.9,0.9,0.9}{
            \makecell[{{p{\linewidth}}}]{
                \texttt{\tiny{[GM$\rangle$P2]}}
                \texttt{Аз съм вашият ръководител на играта и вие играете съвместна игра със следната мрежа като игрално поле:} \\
\\ 
\texttt{```} \\
\texttt{ 1234567} \\
\texttt{╔═╤═════╗} \\
\texttt{║P│S◌D◌◌║ 1} \\
\texttt{╟─┤◌◌◌◌◌║ 2} \\
\texttt{║◌└─┐◌A◌║ 3} \\
\texttt{║◌◌M├─┐◌║ 4} \\
\texttt{║◌T◌│◌│O║ 5} \\
\texttt{╟─┬─┴─┼─╢ 6} \\
\texttt{║◌│◌◌◌│◌║ 7} \\
\texttt{╚═╧═══╧═╝} \\
\\ 
\texttt{```} \\
\\ 
\texttt{* Горният ръб показва x{-}координати, които нарастват надясно, а десният ръб показва y{-}координати, които нарастват надолу.} \\
\texttt{* Следните обекти са разположени на случаен принцип върху вашата мрежа: 'P', 'O', 'T', 'S', 'D', 'A', 'M'.} \\
\\ 
\\ 
\texttt{Другият играч вижда вариант на игралното поле, където обектите са разположени на различни произволни места. Вие не можете да видите полето на другия играч, а той не може да види вашето.} \\
\\ 
\\ 
\texttt{**Цел:**} \\
\\ 
\texttt{И двамата играчи трябва да преместят обектите на своя фон така, че еднакви обекти да се окажат на едни и същи координати. Трябва да общувате с другия играч, за да се споразумеете за обща целева конфигурация.} \\
\\ 
\\ 
\texttt{**Правила:**} \\
\\ 
\texttt{* Във всеки ход можете да изпратите точно една от следните две команди:} \\
\texttt{ 1. `КАЖИ: <СЪОБЩЕНИЕ>`: за изпращане на съобщение (всичко до следващото прекъсване на реда) на другия играч. Аз ще го препратя на вашия партньор.} \\
\texttt{2. `ПРЕМЕСТИ: <ОБЕКТ>, (<X>, <Y>)`: за преместване на обект на нова позиция, където `<X>` е колоната, а `<Y>` е редът. Ще ви информирам дали преместването ви е успешно или не.} \\
\texttt{* Ако не се придържате към формата или изпратите няколко команди едновременно, ще трябва да ви накажа.} \\
\texttt{* Ако и двамата играчи натрупат повече от 16 наказания, и двамата губите играта.} \\
\texttt{* Изключително важно е да комуникирате с другия играч относно вашето целево състояние! Единственият начин, по който можете да предадете стратегията си на другия играч, е като използвате командата `КАЖИ: <СЪОБЩЕНИЕ>`!} \\
\\ 
\\ 
\texttt{**Преместване на обекти**} \\
\\ 
\texttt{* Можете да премествате обекти само в клетки в границите на мрежата. Целевата клетка трябва да е празна, т.е. да съдържа само символа '◌'.} \\
\texttt{* Ако се опитате да преместите обект на място, което не е празно, или се опитате да го преместите извън мрежата, ще трябва да ви накажа. Получавате още един опит.} \\
\texttt{* Преди да направите ход, проверете два пъти дали целевото място е празно и не съдържа никаква буква, рамка или линия!} \\
\\ 
\\ 
\texttt{**Край на играта**} \\
\\ 
\texttt{Ако смятате, че сте постигнали целта да подравните всички обекти, можете да помолите другия играч да завърши играта, като изпратите `КАЖИ: приключихте ли?`. Ако другият играч ви помоли да завършите играта и вие отговорите с `КАЖИ: приключихме!`, играта ще приключи.} \\
\\ 
\texttt{И двамата играчи печелят, ако играта приключи в рамките на 28 рунда, където един рунд се определя като двама играчи, всеки от които изпраща валидна команда.} \\
\\ 
\texttt{**Оценяване:**} \\
\\ 
\texttt{Колкото по{-}близо са еднаквите обекти в двете игрални полета, толкова повече точки получавате. Наказанията намаляват точките ви. Можете ли да подобрите рекорда?} \\
\\ 
\texttt{Другият играч започна играта, като изпрати това съобщение:} \\
\\ 
\texttt{"Здравей! Ето моите обекти и техните позиции: P(1,7), O(1,3), T(3,2), S(3,5), D(5,7), A(3,1), M(6,7). Какви са вашите? Да изберем обща цел!"} \\
\\ 
\texttt{Каква е първата ви команда?} \\
            }
        }
    }
    & \\ \\

    \theutterance \stepcounter{utterance}  
    & & & \multicolumn{4}{p{0.6\linewidth}}{
        \cellcolor[rgb]{0.9,0.9,0.9}{
            \makecell[{{p{\linewidth}}}]{
                \texttt{\tiny{[GM$\langle$P2]}}
                \texttt{Отлично начало от страна на другия играч! Първо, нека споделя моите позиции на обектите, за да можем да започнем да планираме обща цел. Ето моите обекти и техните координати от моето поле:} \\
\\ 
\texttt{{-} P(2,1)} \\
\texttt{{-} S(2,2)} \\
\texttt{{-} D(4,1)} \\
\texttt{{-} A(4,3)} \\
\texttt{{-} M(3,4)} \\
\texttt{{-} T(2,5)} \\
\texttt{{-} O(7,5)} \\
\\ 
\texttt{Моят първи ход ще бъде да споделя тази информация с другия играч, за да започнем да съгласуваме целевите позиции. Ето командата:} \\
\\ 
\texttt{```} \\
\texttt{КАЖИ: Моите обекти са: P(2,1), S(2,2), D(4,1), A(4,3), M(3,4), T(2,5), O(7,5). Да изберем обща цел за всяка двойка!} \\
\texttt{```} \\
            }
        }
    }
     \\ \\

    \theutterance \stepcounter{utterance}  
    & & & \multicolumn{2}{p{0.3\linewidth}}{
        \cellcolor[rgb]{0.9,0.9,0.9}{
            \makecell[{{p{\linewidth}}}]{
                \texttt{\tiny{[GM$|$GM]}}
                \texttt{Отлично начало от страна на другия играч! Първо, нека споделя моите позиции на обектите, за да можем да започнем да планираме обща цел. Ето моите обекти и техните координати от моето поле:} \\
\\ 
\texttt{{-} P(2,1)} \\
\texttt{{-} S(2,2)} \\
\texttt{{-} D(4,1)} \\
\texttt{{-} A(4,3)} \\
\texttt{{-} M(3,4)} \\
\texttt{{-} T(2,5)} \\
\texttt{{-} O(7,5)} \\
\\ 
\texttt{Моят първи ход ще бъде да споделя тази информация с другия играч, за да започнем да съгласуваме целевите позиции. Ето командата:} \\
\\ 
\texttt{```} \\
\texttt{КАЖИ: Моите обекти са: P(2,1), S(2,2), D(4,1), A(4,3), M(3,4), T(2,5), O(7,5). Да изберем обща цел за всяка двойка!} \\
\texttt{```} \\
            }
        }
    }
    & & \\ \\

    \theutterance \stepcounter{utterance}  
    & & & \multicolumn{2}{p{0.3\linewidth}}{
        \cellcolor[rgb]{0.9,0.9,0.9}{
            \makecell[{{p{\linewidth}}}]{
                \texttt{\tiny{[GM$|$GM]}}
                \texttt{Invalid format: head is not empty: 'Отлично начало от страна на другия играч! Първо, нека споделя моите позиции на обектите, за да можем да започнем да планираме обща цел. Ето моите обекти и техните координати от моето поле:} \\
\\ 
\texttt{{-} P(2,1)} \\
\texttt{{-} S(2,2)} \\
\texttt{{-} D(4,1)} \\
\texttt{{-} A(4,3)} \\
\texttt{{-} M(3,4)} \\
\texttt{{-} T(2,5)} \\
\texttt{{-} O(7,5)} \\
\\ 
\texttt{Моят първи ход ще бъде да споделя тази информация с другия играч, за да започнем да съгласуваме целевите позиции. Ето командата:} \\
\\ 
\\ 
\texttt{'} \\
            }
        }
    }
    & & \\ \\

    \theutterance \stepcounter{utterance}  
    & & \multicolumn{4}{p{0.6\linewidth}}{
        \cellcolor[rgb]{0.9,0.9,0.9}{
            \makecell[{{p{\linewidth}}}]{
                \texttt{\tiny{[GM$\rangle$P2]}}
                \texttt{Наказание: Вашето съобщение не трябва да съдържа нищо преди командата!} \\
\\ 
\texttt{Уверете се, че вашият отговор съдържа само `КАЖИ: <СЪОБЩЕНИЕ>` или `ПРЕМЕСТИ: <ОБЕКТ>, (<X>, <Y>)` и нищо друго!} \\
\\ 
\texttt{Колективно сте натрупали 2 от 16 наказания.} \\
\texttt{Моля, опитайте отново!} \\
            }
        }
    }
    & \\ \\

    \theutterance \stepcounter{utterance}  
    & & & \multicolumn{4}{p{0.6\linewidth}}{
        \cellcolor[rgb]{0.9,0.9,0.9}{
            \makecell[{{p{\linewidth}}}]{
                \texttt{\tiny{[GM$\langle$P2]}}
                \texttt{Разбрах! Тук е коригираната команда:} \\
\\ 
\texttt{```} \\
\texttt{КАЖИ: Моите позиции са: P(2,1), S(2,2), D(4,1), A(4,3), M(3,4), T(2,5), O(7,5). Да изберем обща цел за всяка двойка!} \\
\texttt{```} \\
            }
        }
    }
     \\ \\

    \theutterance \stepcounter{utterance}  
    & & & \multicolumn{2}{p{0.3\linewidth}}{
        \cellcolor[rgb]{0.9,0.9,0.9}{
            \makecell[{{p{\linewidth}}}]{
                \texttt{\tiny{[GM$|$GM]}}
                \texttt{Разбрах! Тук е коригираната команда:} \\
\\ 
\texttt{```} \\
\texttt{КАЖИ: Моите позиции са: P(2,1), S(2,2), D(4,1), A(4,3), M(3,4), T(2,5), O(7,5). Да изберем обща цел за всяка двойка!} \\
\texttt{```} \\
            }
        }
    }
    & & \\ \\

    \theutterance \stepcounter{utterance}  
    & & & \multicolumn{2}{p{0.3\linewidth}}{
        \cellcolor[rgb]{0.9,0.9,0.9}{
            \makecell[{{p{\linewidth}}}]{
                \texttt{\tiny{[GM$|$GM]}}
                \texttt{Invalid format: head is not empty: 'Разбрах! Тук е коригираната команда:} \\
\\ 
\\ 
\texttt{'} \\
            }
        }
    }
    & & \\ \\

    \theutterance \stepcounter{utterance}  
    & & \multicolumn{4}{p{0.6\linewidth}}{
        \cellcolor[rgb]{0.9,0.9,0.9}{
            \makecell[{{p{\linewidth}}}]{
                \texttt{\tiny{[GM$\rangle$P2]}}
                \texttt{Наказание: Вашето съобщение не трябва да съдържа нищо преди командата!} \\
\\ 
\texttt{Уверете се, че вашият отговор съдържа само `КАЖИ: <СЪОБЩЕНИЕ>` или `ПРЕМЕСТИ: <ОБЕКТ>, (<X>, <Y>)` и нищо друго!} \\
\\ 
\texttt{Колективно сте натрупали 3 от 16 наказания.} \\
\texttt{Моля, опитайте отново!} \\
            }
        }
    }
    & \\ \\

    \theutterance \stepcounter{utterance}  
    & & & \multicolumn{4}{p{0.6\linewidth}}{
        \cellcolor[rgb]{0.9,0.9,0.9}{
            \makecell[{{p{\linewidth}}}]{
                \texttt{\tiny{[GM$\langle$P2]}}
                \texttt{Ясно, ще спазвам стриктно формата. Ето правилната команда:} \\
\\ 
\texttt{```} \\
\texttt{КАЖИ: P(2,1), S(2,2), D(4,1), A(4,3), M(3,4), T(2,5), O(7,5). Да изберем обща цел.} \\
\texttt{```} \\
            }
        }
    }
     \\ \\

    \theutterance \stepcounter{utterance}  
    & & & \multicolumn{2}{p{0.3\linewidth}}{
        \cellcolor[rgb]{0.9,0.9,0.9}{
            \makecell[{{p{\linewidth}}}]{
                \texttt{\tiny{[GM$|$GM]}}
                \texttt{Ясно, ще спазвам стриктно формата. Ето правилната команда:} \\
\\ 
\texttt{```} \\
\texttt{КАЖИ: P(2,1), S(2,2), D(4,1), A(4,3), M(3,4), T(2,5), O(7,5). Да изберем обща цел.} \\
\texttt{```} \\
            }
        }
    }
    & & \\ \\

    \theutterance \stepcounter{utterance}  
    & & & \multicolumn{2}{p{0.3\linewidth}}{
        \cellcolor[rgb]{0.9,0.9,0.9}{
            \makecell[{{p{\linewidth}}}]{
                \texttt{\tiny{[GM$|$GM]}}
                \texttt{Invalid format: head is not empty: 'Ясно, ще спазвам стриктно формата. Ето правилната команда:} \\
\\ 
\\ 
\texttt{'} \\
            }
        }
    }
    & & \\ \\

    \theutterance \stepcounter{utterance}  
    & & \multicolumn{4}{p{0.6\linewidth}}{
        \cellcolor[rgb]{0.9,0.9,0.9}{
            \makecell[{{p{\linewidth}}}]{
                \texttt{\tiny{[GM$\rangle$P2]}}
                \texttt{Наказание: Вашето съобщение не трябва да съдържа нищо преди командата!} \\
\\ 
\texttt{Уверете се, че вашият отговор съдържа само `КАЖИ: <СЪОБЩЕНИЕ>` или `ПРЕМЕСТИ: <ОБЕКТ>, (<X>, <Y>)` и нищо друго!} \\
\\ 
\texttt{Колективно сте натрупали 4 от 16 наказания.} \\
\texttt{Моля, опитайте отново!} \\
            }
        }
    }
    & \\ \\

    \theutterance \stepcounter{utterance}  
    & & & \multicolumn{4}{p{0.6\linewidth}}{
        \cellcolor[rgb]{0.9,0.9,0.9}{
            \makecell[{{p{\linewidth}}}]{
                \texttt{\tiny{[GM$\langle$P2]}}
                \texttt{Съжалявам за грешките. Ето правилната команда:} \\
\\ 
\texttt{```} \\
\texttt{КАЖИ: Ето моите позиции: P(2,1) S(2,2) D(4,1) A(4,3) M(3,4) T(2,5) O(7,5)} \\
\texttt{```} \\
            }
        }
    }
     \\ \\

    \theutterance \stepcounter{utterance}  
    & & & \multicolumn{2}{p{0.3\linewidth}}{
        \cellcolor[rgb]{0.9,0.9,0.9}{
            \makecell[{{p{\linewidth}}}]{
                \texttt{\tiny{[GM$|$GM]}}
                \texttt{Съжалявам за грешките. Ето правилната команда:} \\
\\ 
\texttt{```} \\
\texttt{КАЖИ: Ето моите позиции: P(2,1) S(2,2) D(4,1) A(4,3) M(3,4) T(2,5) O(7,5)} \\
\texttt{```} \\
            }
        }
    }
    & & \\ \\

    \theutterance \stepcounter{utterance}  
    & & & \multicolumn{2}{p{0.3\linewidth}}{
        \cellcolor[rgb]{0.9,0.9,0.9}{
            \makecell[{{p{\linewidth}}}]{
                \texttt{\tiny{[GM$|$GM]}}
                \texttt{Invalid format: head is not empty: 'Съжалявам за грешките. Ето правилната команда:} \\
\\ 
\\ 
\texttt{'} \\
            }
        }
    }
    & & \\ \\

    \theutterance \stepcounter{utterance}  
    & & \multicolumn{4}{p{0.6\linewidth}}{
        \cellcolor[rgb]{0.9,0.9,0.9}{
            \makecell[{{p{\linewidth}}}]{
                \texttt{\tiny{[GM$\rangle$P2]}}
                \texttt{Наказание: Вашето съобщение не трябва да съдържа нищо преди командата!} \\
\\ 
\texttt{Уверете се, че вашият отговор съдържа само `КАЖИ: <СЪОБЩЕНИЕ>` или `ПРЕМЕСТИ: <ОБЕКТ>, (<X>, <Y>)` и нищо друго!} \\
\\ 
\texttt{Колективно сте натрупали 5 от 16 наказания.} \\
\texttt{Моля, опитайте отново!} \\
            }
        }
    }
    & \\ \\

    \theutterance \stepcounter{utterance}  
    & & & \multicolumn{4}{p{0.6\linewidth}}{
        \cellcolor[rgb]{0.9,0.9,0.9}{
            \makecell[{{p{\linewidth}}}]{
                \texttt{\tiny{[GM$\langle$P2]}}
                \texttt{Разбрах {-} ще спазвам формата стриктно. Ето чистата команда без никакви допълнения:} \\
\\ 
\texttt{```} \\
\texttt{КАЖИ: P(2,1) S(2,2) D(4,1) A(4,3) M(3,4) T(2,5) O(7,5)} \\
\texttt{```} \\
            }
        }
    }
     \\ \\

    \theutterance \stepcounter{utterance}  
    & & & \multicolumn{2}{p{0.3\linewidth}}{
        \cellcolor[rgb]{0.9,0.9,0.9}{
            \makecell[{{p{\linewidth}}}]{
                \texttt{\tiny{[GM$|$GM]}}
                \texttt{Разбрах {-} ще спазвам формата стриктно. Ето чистата команда без никакви допълнения:} \\
\\ 
\texttt{```} \\
\texttt{КАЖИ: P(2,1) S(2,2) D(4,1) A(4,3) M(3,4) T(2,5) O(7,5)} \\
\texttt{```} \\
            }
        }
    }
    & & \\ \\

    \theutterance \stepcounter{utterance}  
    & & & \multicolumn{2}{p{0.3\linewidth}}{
        \cellcolor[rgb]{0.9,0.9,0.9}{
            \makecell[{{p{\linewidth}}}]{
                \texttt{\tiny{[GM$|$GM]}}
                \texttt{Invalid format: head is not empty: 'Разбрах {-} ще спазвам формата стриктно. Ето чистата команда без никакви допълнения:} \\
\\ 
\\ 
\texttt{'} \\
            }
        }
    }
    & & \\ \\

    \theutterance \stepcounter{utterance}  
    & & \multicolumn{4}{p{0.6\linewidth}}{
        \cellcolor[rgb]{0.9,0.9,0.9}{
            \makecell[{{p{\linewidth}}}]{
                \texttt{\tiny{[GM$\rangle$P2]}}
                \texttt{Наказание: Вашето съобщение не трябва да съдържа нищо преди командата!} \\
\\ 
\texttt{Уверете се, че вашият отговор съдържа само `КАЖИ: <СЪОБЩЕНИЕ>` или `ПРЕМЕСТИ: <ОБЕКТ>, (<X>, <Y>)` и нищо друго!} \\
\\ 
\texttt{Колективно сте натрупали 6 от 16 наказания.} \\
\texttt{Моля, опитайте отново!} \\
            }
        }
    }
    & \\ \\

    \theutterance \stepcounter{utterance}  
    & & & \multicolumn{4}{p{0.6\linewidth}}{
        \cellcolor[rgb]{0.9,0.9,0.9}{
            \makecell[{{p{\linewidth}}}]{
                \texttt{\tiny{[GM$\langle$P2]}}
                \texttt{Ясно, ще изпратя само чистата информация без никакви допълнителни думи:} \\
\\ 
\texttt{```} \\
\texttt{КАЖИ: P(2,1),S(2,2),D(4,1),A(4,3),M(3,4),T(2,5),O(7,5)} \\
\texttt{```} \\
            }
        }
    }
     \\ \\

    \theutterance \stepcounter{utterance}  
    & & & \multicolumn{2}{p{0.3\linewidth}}{
        \cellcolor[rgb]{0.9,0.9,0.9}{
            \makecell[{{p{\linewidth}}}]{
                \texttt{\tiny{[GM$|$GM]}}
                \texttt{Ясно, ще изпратя само чистата информация без никакви допълнителни думи:} \\
\\ 
\texttt{```} \\
\texttt{КАЖИ: P(2,1),S(2,2),D(4,1),A(4,3),M(3,4),T(2,5),O(7,5)} \\
\texttt{```} \\
            }
        }
    }
    & & \\ \\

    \theutterance \stepcounter{utterance}  
    & & & \multicolumn{2}{p{0.3\linewidth}}{
        \cellcolor[rgb]{0.9,0.9,0.9}{
            \makecell[{{p{\linewidth}}}]{
                \texttt{\tiny{[GM$|$GM]}}
                \texttt{Invalid format: head is not empty: 'Ясно, ще изпратя само чистата информация без никакви допълнителни думи:} \\
\\ 
\\ 
\texttt{'} \\
            }
        }
    }
    & & \\ \\

    \theutterance \stepcounter{utterance}  
    & & \multicolumn{4}{p{0.6\linewidth}}{
        \cellcolor[rgb]{0.9,0.9,0.9}{
            \makecell[{{p{\linewidth}}}]{
                \texttt{\tiny{[GM$\rangle$P2]}}
                \texttt{Наказание: Вашето съобщение не трябва да съдържа нищо преди командата!} \\
\\ 
\texttt{Уверете се, че вашият отговор съдържа само `КАЖИ: <СЪОБЩЕНИЕ>` или `ПРЕМЕСТИ: <ОБЕКТ>, (<X>, <Y>)` и нищо друго!} \\
\\ 
\texttt{Колективно сте натрупали 7 от 16 наказания.} \\
\texttt{Моля, опитайте отново!} \\
            }
        }
    }
    & \\ \\

    \theutterance \stepcounter{utterance}  
    & & & \multicolumn{4}{p{0.6\linewidth}}{
        \cellcolor[rgb]{0.9,0.9,0.9}{
            \makecell[{{p{\linewidth}}}]{
                \texttt{\tiny{[GM$\langle$P2]}}
                \texttt{Съжалявам за грешките. Ето абсолютно минималната команда:} \\
\\ 
\texttt{```} \\
\texttt{КАЖИ: P2,1 S2,2 D4,1 A4,3 M3,4 T2,5 O7,5} \\
\texttt{```} \\
            }
        }
    }
     \\ \\

    \theutterance \stepcounter{utterance}  
    & & & \multicolumn{2}{p{0.3\linewidth}}{
        \cellcolor[rgb]{0.9,0.9,0.9}{
            \makecell[{{p{\linewidth}}}]{
                \texttt{\tiny{[GM$|$GM]}}
                \texttt{Съжалявам за грешките. Ето абсолютно минималната команда:} \\
\\ 
\texttt{```} \\
\texttt{КАЖИ: P2,1 S2,2 D4,1 A4,3 M3,4 T2,5 O7,5} \\
\texttt{```} \\
            }
        }
    }
    & & \\ \\

    \theutterance \stepcounter{utterance}  
    & & & \multicolumn{2}{p{0.3\linewidth}}{
        \cellcolor[rgb]{0.9,0.9,0.9}{
            \makecell[{{p{\linewidth}}}]{
                \texttt{\tiny{[GM$|$GM]}}
                \texttt{Invalid format: head is not empty: 'Съжалявам за грешките. Ето абсолютно минималната команда:} \\
\\ 
\\ 
\texttt{'} \\
            }
        }
    }
    & & \\ \\

    \theutterance \stepcounter{utterance}  
    & & \multicolumn{4}{p{0.6\linewidth}}{
        \cellcolor[rgb]{0.9,0.9,0.9}{
            \makecell[{{p{\linewidth}}}]{
                \texttt{\tiny{[GM$\rangle$P2]}}
                \texttt{Наказание: Вашето съобщение не трябва да съдържа нищо преди командата!} \\
\\ 
\texttt{Уверете се, че вашият отговор съдържа само `КАЖИ: <СЪОБЩЕНИЕ>` или `ПРЕМЕСТИ: <ОБЕКТ>, (<X>, <Y>)` и нищо друго!} \\
\\ 
\texttt{Колективно сте натрупали 8 от 16 наказания.} \\
\texttt{Моля, опитайте отново!} \\
            }
        }
    }
    & \\ \\

    \theutterance \stepcounter{utterance}  
    & & & \multicolumn{4}{p{0.6\linewidth}}{
        \cellcolor[rgb]{0.9,0.9,0.9}{
            \makecell[{{p{\linewidth}}}]{
                \texttt{\tiny{[GM$\langle$P2]}}
                \texttt{Ясно {-} ще изпратя само координатите в най{-}простия възможен формат:} \\
\\ 
\texttt{```} \\
\texttt{КАЖИ: P(2,1)S(2,2)D(4,1)A(4,3)M(3,4)T(2,5)O(7,5)} \\
\texttt{```} \\
            }
        }
    }
     \\ \\

    \theutterance \stepcounter{utterance}  
    & & & \multicolumn{2}{p{0.3\linewidth}}{
        \cellcolor[rgb]{0.9,0.9,0.9}{
            \makecell[{{p{\linewidth}}}]{
                \texttt{\tiny{[GM$|$GM]}}
                \texttt{Ясно {-} ще изпратя само координатите в най{-}простия възможен формат:} \\
\\ 
\texttt{```} \\
\texttt{КАЖИ: P(2,1)S(2,2)D(4,1)A(4,3)M(3,4)T(2,5)O(7,5)} \\
\texttt{```} \\
            }
        }
    }
    & & \\ \\

    \theutterance \stepcounter{utterance}  
    & & & \multicolumn{2}{p{0.3\linewidth}}{
        \cellcolor[rgb]{0.9,0.9,0.9}{
            \makecell[{{p{\linewidth}}}]{
                \texttt{\tiny{[GM$|$GM]}}
                \texttt{Invalid format: head is not empty: 'Ясно {-} ще изпратя само координатите в най{-}простия възможен формат:} \\
\\ 
\\ 
\texttt{'} \\
            }
        }
    }
    & & \\ \\

    \theutterance \stepcounter{utterance}  
    & & \multicolumn{4}{p{0.6\linewidth}}{
        \cellcolor[rgb]{0.9,0.9,0.9}{
            \makecell[{{p{\linewidth}}}]{
                \texttt{\tiny{[GM$\rangle$P2]}}
                \texttt{Наказание: Вашето съобщение не трябва да съдържа нищо преди командата!} \\
\\ 
\texttt{Уверете се, че вашият отговор съдържа само `КАЖИ: <СЪОБЩЕНИЕ>` или `ПРЕМЕСТИ: <ОБЕКТ>, (<X>, <Y>)` и нищо друго!} \\
\\ 
\texttt{Колективно сте натрупали 9 от 16 наказания.} \\
\texttt{Моля, опитайте отново!} \\
            }
        }
    }
    & \\ \\

    \theutterance \stepcounter{utterance}  
    & & & \multicolumn{4}{p{0.6\linewidth}}{
        \cellcolor[rgb]{0.9,0.9,0.9}{
            \makecell[{{p{\linewidth}}}]{
                \texttt{\tiny{[GM$\langle$P2]}}
                \texttt{Моля, приемете следната команда без никакви допълнения:} \\
\\ 
\texttt{```} \\
\texttt{КАЖИ: (2,1)P(2,2)S(4,1)D(4,3)A(3,4)M(2,5)T(7,5)O} \\
\texttt{```} \\
            }
        }
    }
     \\ \\

    \theutterance \stepcounter{utterance}  
    & & & \multicolumn{2}{p{0.3\linewidth}}{
        \cellcolor[rgb]{0.9,0.9,0.9}{
            \makecell[{{p{\linewidth}}}]{
                \texttt{\tiny{[GM$|$GM]}}
                \texttt{Моля, приемете следната команда без никакви допълнения:} \\
\\ 
\texttt{```} \\
\texttt{КАЖИ: (2,1)P(2,2)S(4,1)D(4,3)A(3,4)M(2,5)T(7,5)O} \\
\texttt{```} \\
            }
        }
    }
    & & \\ \\

    \theutterance \stepcounter{utterance}  
    & & & \multicolumn{2}{p{0.3\linewidth}}{
        \cellcolor[rgb]{0.9,0.9,0.9}{
            \makecell[{{p{\linewidth}}}]{
                \texttt{\tiny{[GM$|$GM]}}
                \texttt{Invalid format: head is not empty: 'Моля, приемете следната команда без никакви допълнения:} \\
\\ 
\\ 
\texttt{'} \\
            }
        }
    }
    & & \\ \\

    \theutterance \stepcounter{utterance}  
    & & \multicolumn{4}{p{0.6\linewidth}}{
        \cellcolor[rgb]{0.9,0.9,0.9}{
            \makecell[{{p{\linewidth}}}]{
                \texttt{\tiny{[GM$\rangle$P2]}}
                \texttt{Наказание: Вашето съобщение не трябва да съдържа нищо преди командата!} \\
\\ 
\texttt{Уверете се, че вашият отговор съдържа само `КАЖИ: <СЪОБЩЕНИЕ>` или `ПРЕМЕСТИ: <ОБЕКТ>, (<X>, <Y>)` и нищо друго!} \\
\\ 
\texttt{Колективно сте натрупали 10 от 16 наказания.} \\
\texttt{Моля, опитайте отново!} \\
            }
        }
    }
    & \\ \\

    \theutterance \stepcounter{utterance}  
    & & & \multicolumn{4}{p{0.6\linewidth}}{
        \cellcolor[rgb]{0.9,0.9,0.9}{
            \makecell[{{p{\linewidth}}}]{
                \texttt{\tiny{[GM$\langle$P2]}}
                \texttt{Разбрах проблема {-} ще изпратя само чистата информация във формат, който лесно може да се декодира:} \\
\\ 
\texttt{```} \\
\texttt{КАЖИ: P=2,1;S=2,2;D=4,1;A=4,3;M=3,4;T=2,5;O=7,5} \\
\texttt{```} \\
\\ 
\texttt{Ако и това не е приемливо, моля да ми кажете какъв точно формат се очаква, за да избегна бъдещи наказания. Благодаря!} \\
            }
        }
    }
     \\ \\

    \theutterance \stepcounter{utterance}  
    & & & \multicolumn{2}{p{0.3\linewidth}}{
        \cellcolor[rgb]{0.9,0.9,0.9}{
            \makecell[{{p{\linewidth}}}]{
                \texttt{\tiny{[GM$|$GM]}}
                \texttt{Разбрах проблема {-} ще изпратя само чистата информация във формат, който лесно може да се декодира:} \\
\\ 
\texttt{```} \\
\texttt{КАЖИ: P=2,1;S=2,2;D=4,1;A=4,3;M=3,4;T=2,5;O=7,5} \\
\texttt{```} \\
\\ 
\texttt{Ако и това не е приемливо, моля да ми кажете какъв точно формат се очаква, за да избегна бъдещи наказания. Благодаря!} \\
            }
        }
    }
    & & \\ \\

    \theutterance \stepcounter{utterance}  
    & & & \multicolumn{2}{p{0.3\linewidth}}{
        \cellcolor[rgb]{0.9,0.9,0.9}{
            \makecell[{{p{\linewidth}}}]{
                \texttt{\tiny{[GM$|$GM]}}
                \texttt{Invalid format: head and tail are not empty} \\
\texttt{head: 'Разбрах проблема {-} ще изпратя само чистата информация във формат, който лесно може да се декодира:} \\
\\ 
\\ 
\texttt{'} \\
\texttt{tail: '} \\
\\ 
\\ 
\texttt{Ако и това не е приемливо, моля да ми кажете какъв точно формат се очаква, за да избегна бъдещи наказания. Благодаря!'} \\
            }
        }
    }
    & & \\ \\

    \theutterance \stepcounter{utterance}  
    & & \multicolumn{4}{p{0.6\linewidth}}{
        \cellcolor[rgb]{0.9,0.9,0.9}{
            \makecell[{{p{\linewidth}}}]{
                \texttt{\tiny{[GM$\rangle$P2]}}
                \texttt{Наказание: Вашето съобщение не трябва да съдържа нищо преди или след командата!} \\
\\ 
\texttt{Уверете се, че вашият отговор съдържа само `КАЖИ: <СЪОБЩЕНИЕ>` или `ПРЕМЕСТИ: <ОБЕКТ>, (<X>, <Y>)` и нищо друго!} \\
\\ 
\texttt{Колективно сте натрупали 11 от 16 наказания.} \\
\texttt{Моля, опитайте отново!} \\
            }
        }
    }
    & \\ \\

    \theutterance \stepcounter{utterance}  
    & & & \multicolumn{4}{p{0.6\linewidth}}{
        \cellcolor[rgb]{0.9,0.9,0.9}{
            \makecell[{{p{\linewidth}}}]{
                \texttt{\tiny{[GM$\langle$P2]}}
                \texttt{Съжалявам за затрудненията. Ето перфектно форматираната команда:} \\
\\ 
\texttt{```} \\
\texttt{КАЖИ: Моите обекти: P на 2,1; S на 2,2; D на 4,1; A на 4,3; M на 3,4; T на 2,5; O на 7,5} \\
\texttt{```} \\
\\ 
\texttt{Сега ще изпратя само чистата информация без каквито и да било допълнителни думи:} \\
\\ 
\texttt{```} \\
\texttt{КАЖИ: P2,1 S2,2 D4,1 A4,3 M3,4 T2,5 O7,5} \\
\texttt{```} \\
            }
        }
    }
     \\ \\

    \theutterance \stepcounter{utterance}  
    & & & \multicolumn{2}{p{0.3\linewidth}}{
        \cellcolor[rgb]{0.9,0.9,0.9}{
            \makecell[{{p{\linewidth}}}]{
                \texttt{\tiny{[GM$|$GM]}}
                \texttt{Съжалявам за затрудненията. Ето перфектно форматираната команда:} \\
\\ 
\texttt{```} \\
\texttt{КАЖИ: Моите обекти: P на 2,1; S на 2,2; D на 4,1; A на 4,3; M на 3,4; T на 2,5; O на 7,5} \\
\texttt{```} \\
\\ 
\texttt{Сега ще изпратя само чистата информация без каквито и да било допълнителни думи:} \\
\\ 
\texttt{```} \\
\texttt{КАЖИ: P2,1 S2,2 D4,1 A4,3 M3,4 T2,5 O7,5} \\
\texttt{```} \\
            }
        }
    }
    & & \\ \\

    \theutterance \stepcounter{utterance}  
    & & & \multicolumn{2}{p{0.3\linewidth}}{
        \cellcolor[rgb]{0.9,0.9,0.9}{
            \makecell[{{p{\linewidth}}}]{
                \texttt{\tiny{[GM$|$GM]}}
                \texttt{Invalid format: head is not empty: 'Съжалявам за затрудненията. Ето перфектно форматираната команда:} \\
\\ 
\\ 
\texttt{КАЖИ: Моите обекти: P на 2,1; S на 2,2; D на 4,1; A на 4,3; M на 3,4; T на 2,5; O на 7,5} \\
\\ 
\\ 
\texttt{Сега ще изпратя само чистата информация без каквито и да било допълнителни думи:} \\
\\ 
\\ 
\texttt{'} \\
            }
        }
    }
    & & \\ \\

    \theutterance \stepcounter{utterance}  
    & & \multicolumn{4}{p{0.6\linewidth}}{
        \cellcolor[rgb]{0.9,0.9,0.9}{
            \makecell[{{p{\linewidth}}}]{
                \texttt{\tiny{[GM$\rangle$P2]}}
                \texttt{Наказание: Вашето съобщение не трябва да съдържа нищо преди командата!} \\
\\ 
\texttt{Уверете се, че вашият отговор съдържа само `КАЖИ: <СЪОБЩЕНИЕ>` или `ПРЕМЕСТИ: <ОБЕКТ>, (<X>, <Y>)` и нищо друго!} \\
\\ 
\texttt{Колективно сте натрупали 12 от 16 наказания.} \\
\texttt{Моля, опитайте отново!} \\
            }
        }
    }
    & \\ \\

    \theutterance \stepcounter{utterance}  
    & & & \multicolumn{4}{p{0.6\linewidth}}{
        \cellcolor[rgb]{0.9,0.9,0.9}{
            \makecell[{{p{\linewidth}}}]{
                \texttt{\tiny{[GM$\langle$P2]}}
                \texttt{Ясно {-} ще изпратя само координатите в най{-}чист вид без никакви букви или символи:} \\
\\ 
\texttt{```} \\
\texttt{КАЖИ: 2,1 2,2 4,1 4,3 3,4 2,5 7,5} \\
\texttt{```} \\
\\ 
\texttt{Ако това все още не е приемливо, моля да ми кажете точния формат, който се очаква за съобщението. Благодаря!} \\
            }
        }
    }
     \\ \\

    \theutterance \stepcounter{utterance}  
    & & & \multicolumn{2}{p{0.3\linewidth}}{
        \cellcolor[rgb]{0.9,0.9,0.9}{
            \makecell[{{p{\linewidth}}}]{
                \texttt{\tiny{[GM$|$GM]}}
                \texttt{Ясно {-} ще изпратя само координатите в най{-}чист вид без никакви букви или символи:} \\
\\ 
\texttt{```} \\
\texttt{КАЖИ: 2,1 2,2 4,1 4,3 3,4 2,5 7,5} \\
\texttt{```} \\
\\ 
\texttt{Ако това все още не е приемливо, моля да ми кажете точния формат, който се очаква за съобщението. Благодаря!} \\
            }
        }
    }
    & & \\ \\

    \theutterance \stepcounter{utterance}  
    & & & \multicolumn{2}{p{0.3\linewidth}}{
        \cellcolor[rgb]{0.9,0.9,0.9}{
            \makecell[{{p{\linewidth}}}]{
                \texttt{\tiny{[GM$|$GM]}}
                \texttt{Invalid format: head and tail are not empty} \\
\texttt{head: 'Ясно {-} ще изпратя само координатите в най{-}чист вид без никакви букви или символи:} \\
\\ 
\\ 
\texttt{'} \\
\texttt{tail: '} \\
\\ 
\\ 
\texttt{Ако това все още не е приемливо, моля да ми кажете точния формат, който се очаква за съобщението. Благодаря!'} \\
            }
        }
    }
    & & \\ \\

    \theutterance \stepcounter{utterance}  
    & & \multicolumn{4}{p{0.6\linewidth}}{
        \cellcolor[rgb]{0.9,0.9,0.9}{
            \makecell[{{p{\linewidth}}}]{
                \texttt{\tiny{[GM$\rangle$P2]}}
                \texttt{Наказание: Вашето съобщение не трябва да съдържа нищо преди или след командата!} \\
\\ 
\texttt{Уверете се, че вашият отговор съдържа само `КАЖИ: <СЪОБЩЕНИЕ>` или `ПРЕМЕСТИ: <ОБЕКТ>, (<X>, <Y>)` и нищо друго!} \\
\\ 
\texttt{Колективно сте натрупали 13 от 16 наказания.} \\
\texttt{Моля, опитайте отново!} \\
            }
        }
    }
    & \\ \\

    \theutterance \stepcounter{utterance}  
    & & & \multicolumn{4}{p{0.6\linewidth}}{
        \cellcolor[rgb]{0.9,0.9,0.9}{
            \makecell[{{p{\linewidth}}}]{
                \texttt{\tiny{[GM$\langle$P2]}}
                \texttt{Разбирам проблема {-} ще изпратя само координатите на обектите в строго определения формат:} \\
\\ 
\texttt{```} \\
\texttt{КАЖИ: P(2,1) O(7,5) T(2,5) S(2,2) D(4,1) A(4,3) M(3,4)} \\
\texttt{```} \\
\\ 
\texttt{Това е последната ми опция за предаване на информацията в правилния формат. Ако и това не е приемливо, моля да ми посочите точния формат, който се изисква.} \\
            }
        }
    }
     \\ \\

    \theutterance \stepcounter{utterance}  
    & & & \multicolumn{2}{p{0.3\linewidth}}{
        \cellcolor[rgb]{0.9,0.9,0.9}{
            \makecell[{{p{\linewidth}}}]{
                \texttt{\tiny{[GM$|$GM]}}
                \texttt{Разбирам проблема {-} ще изпратя само координатите на обектите в строго определения формат:} \\
\\ 
\texttt{```} \\
\texttt{КАЖИ: P(2,1) O(7,5) T(2,5) S(2,2) D(4,1) A(4,3) M(3,4)} \\
\texttt{```} \\
\\ 
\texttt{Това е последната ми опция за предаване на информацията в правилния формат. Ако и това не е приемливо, моля да ми посочите точния формат, който се изисква.} \\
            }
        }
    }
    & & \\ \\

    \theutterance \stepcounter{utterance}  
    & & & \multicolumn{2}{p{0.3\linewidth}}{
        \cellcolor[rgb]{0.9,0.9,0.9}{
            \makecell[{{p{\linewidth}}}]{
                \texttt{\tiny{[GM$|$GM]}}
                \texttt{Invalid format: head and tail are not empty} \\
\texttt{head: 'Разбирам проблема {-} ще изпратя само координатите на обектите в строго определения формат:} \\
\\ 
\\ 
\texttt{'} \\
\texttt{tail: '} \\
\\ 
\\ 
\texttt{Това е последната ми опция за предаване на информацията в правилния формат. Ако и това не е приемливо, моля да ми посочите точния формат, който се изисква.'} \\
            }
        }
    }
    & & \\ \\

    \theutterance \stepcounter{utterance}  
    & & \multicolumn{4}{p{0.6\linewidth}}{
        \cellcolor[rgb]{0.9,0.9,0.9}{
            \makecell[{{p{\linewidth}}}]{
                \texttt{\tiny{[GM$\rangle$P2]}}
                \texttt{Наказание: Вашето съобщение не трябва да съдържа нищо преди или след командата!} \\
\\ 
\texttt{Уверете се, че вашият отговор съдържа само `КАЖИ: <СЪОБЩЕНИЕ>` или `ПРЕМЕСТИ: <ОБЕКТ>, (<X>, <Y>)` и нищо друго!} \\
\\ 
\texttt{Колективно сте натрупали 14 от 16 наказания.} \\
\texttt{Моля, опитайте отново!} \\
            }
        }
    }
    & \\ \\

    \theutterance \stepcounter{utterance}  
    & & & \multicolumn{4}{p{0.6\linewidth}}{
        \cellcolor[rgb]{0.9,0.9,0.9}{
            \makecell[{{p{\linewidth}}}]{
                \texttt{\tiny{[GM$\langle$P2]}}
                \texttt{Съжалявам за многократните грешки. Ето абсолютно минималната и правилна команда:} \\
\\ 
\texttt{```} \\
\texttt{КАЖИ: P(2,1)S(2,2)D(4,1)A(4,3)M(3,4)T(2,5)O(7,5)} \\
\texttt{```} \\
\\ 
\texttt{Ако това все още не е приемливо, моля да ми посочите точния формат, който трябва да следвам. Благодаря за търпението!} \\
            }
        }
    }
     \\ \\

    \theutterance \stepcounter{utterance}  
    & & & \multicolumn{2}{p{0.3\linewidth}}{
        \cellcolor[rgb]{0.9,0.9,0.9}{
            \makecell[{{p{\linewidth}}}]{
                \texttt{\tiny{[GM$|$GM]}}
                \texttt{Съжалявам за многократните грешки. Ето абсолютно минималната и правилна команда:} \\
\\ 
\texttt{```} \\
\texttt{КАЖИ: P(2,1)S(2,2)D(4,1)A(4,3)M(3,4)T(2,5)O(7,5)} \\
\texttt{```} \\
\\ 
\texttt{Ако това все още не е приемливо, моля да ми посочите точния формат, който трябва да следвам. Благодаря за търпението!} \\
            }
        }
    }
    & & \\ \\

    \theutterance \stepcounter{utterance}  
    & & & \multicolumn{2}{p{0.3\linewidth}}{
        \cellcolor[rgb]{0.9,0.9,0.9}{
            \makecell[{{p{\linewidth}}}]{
                \texttt{\tiny{[GM$|$GM]}}
                \texttt{Invalid format: head and tail are not empty} \\
\texttt{head: 'Съжалявам за многократните грешки. Ето абсолютно минималната и правилна команда:} \\
\\ 
\\ 
\texttt{'} \\
\texttt{tail: '} \\
\\ 
\\ 
\texttt{Ако това все още не е приемливо, моля да ми посочите точния формат, който трябва да следвам. Благодаря за търпението!'} \\
            }
        }
    }
    & & \\ \\

    \theutterance \stepcounter{utterance}  
    & & \multicolumn{4}{p{0.6\linewidth}}{
        \cellcolor[rgb]{0.9,0.9,0.9}{
            \makecell[{{p{\linewidth}}}]{
                \texttt{\tiny{[GM$\rangle$P2]}}
                \texttt{Наказание: Вашето съобщение не трябва да съдържа нищо преди или след командата!} \\
\\ 
\texttt{Уверете се, че вашият отговор съдържа само `КАЖИ: <СЪОБЩЕНИЕ>` или `ПРЕМЕСТИ: <ОБЕКТ>, (<X>, <Y>)` и нищо друго!} \\
\\ 
\texttt{Колективно сте натрупали 15 от 16 наказания.} \\
\texttt{Моля, опитайте отново!} \\
            }
        }
    }
    & \\ \\

    \theutterance \stepcounter{utterance}  
    & & & \multicolumn{4}{p{0.6\linewidth}}{
        \cellcolor[rgb]{0.9,0.9,0.9}{
            \makecell[{{p{\linewidth}}}]{
                \texttt{\tiny{[GM$\langle$P2]}}
                \texttt{Разбрах проблема {-} ще изпратя само една координата като пример, за да избегна объркване:} \\
\\ 
\texttt{```} \\
\texttt{КАЖИ: P(2,1)} \\
\texttt{```} \\
\\ 
\texttt{След това можем да продължим с останалите обекти в отделни съобщения, ако това е приемливо. Ако не, моля да ми посочите точния формат, който трябва да следвам за цялата информация.} \\
            }
        }
    }
     \\ \\

    \theutterance \stepcounter{utterance}  
    & & & \multicolumn{2}{p{0.3\linewidth}}{
        \cellcolor[rgb]{0.9,0.9,0.9}{
            \makecell[{{p{\linewidth}}}]{
                \texttt{\tiny{[GM$|$GM]}}
                \texttt{Разбрах проблема {-} ще изпратя само една координата като пример, за да избегна объркване:} \\
\\ 
\texttt{```} \\
\texttt{КАЖИ: P(2,1)} \\
\texttt{```} \\
\\ 
\texttt{След това можем да продължим с останалите обекти в отделни съобщения, ако това е приемливо. Ако не, моля да ми посочите точния формат, който трябва да следвам за цялата информация.} \\
            }
        }
    }
    & & \\ \\

    \theutterance \stepcounter{utterance}  
    & & & \multicolumn{2}{p{0.3\linewidth}}{
        \cellcolor[rgb]{0.9,0.9,0.9}{
            \makecell[{{p{\linewidth}}}]{
                \texttt{\tiny{[GM$|$GM]}}
                \texttt{Invalid format: head and tail are not empty} \\
\texttt{head: 'Разбрах проблема {-} ще изпратя само една координата като пример, за да избегна объркване:} \\
\\ 
\\ 
\texttt{'} \\
\texttt{tail: '} \\
\\ 
\\ 
\texttt{След това можем да продължим с останалите обекти в отделни съобщения, ако това е приемливо. Ако не, моля да ми посочите точния формат, който трябва да следвам за цялата информация.'} \\
            }
        }
    }
    & & \\ \\

    \theutterance \stepcounter{utterance}  
    & & \multicolumn{4}{p{0.6\linewidth}}{
        \cellcolor[rgb]{0.9,0.9,0.9}{
            \makecell[{{p{\linewidth}}}]{
                \texttt{\tiny{[GM$\rangle$P2]}}
                \texttt{Наказание: Вашето съобщение не трябва да съдържа нищо преди или след командата!} \\
\\ 
\texttt{Уверете се, че вашият отговор съдържа само `КАЖИ: <СЪОБЩЕНИЕ>` или `ПРЕМЕСТИ: <ОБЕКТ>, (<X>, <Y>)` и нищо друго!} \\
\\ 
\texttt{Колективно сте натрупали 16 от 16 наказания.} \\
\texttt{Моля, опитайте отново!} \\
            }
        }
    }
    & \\ \\

    \theutterance \stepcounter{utterance}  
    & & & \multicolumn{4}{p{0.6\linewidth}}{
        \cellcolor[rgb]{0.9,0.9,0.9}{
            \makecell[{{p{\linewidth}}}]{
                \texttt{\tiny{[GM$\langle$P2]}}
                \texttt{Съжалявам за всички грешки. Ето перфектно форматираната команда:} \\
\\ 
\texttt{**ПРЕМЕСТИ: P, (3,1)**} \\
\\ 
\texttt{Това е валидна команда за преместване, която ще ми помогне да избегна наказания, докато планираме следващите стъпки. Моля потвърдете дали е приета. След това ще продължа с комуникацията за координиране на целите.} \\
            }
        }
    }
     \\ \\

    \theutterance \stepcounter{utterance}  
    & & & \multicolumn{2}{p{0.3\linewidth}}{
        \cellcolor[rgb]{0.9,0.9,0.9}{
            \makecell[{{p{\linewidth}}}]{
                \texttt{\tiny{[GM$|$GM]}}
                \texttt{Съжалявам за всички грешки. Ето перфектно форматираната команда:} \\
\\ 
\texttt{**ПРЕМЕСТИ: P, (3,1)**} \\
\\ 
\texttt{Това е валидна команда за преместване, която ще ми помогне да избегна наказания, докато планираме следващите стъпки. Моля потвърдете дали е приета. След това ще продължа с комуникацията за координиране на целите.} \\
            }
        }
    }
    & & \\ \\

    \theutterance \stepcounter{utterance}  
    & & & \multicolumn{2}{p{0.3\linewidth}}{
        \cellcolor[rgb]{0.9,0.9,0.9}{
            \makecell[{{p{\linewidth}}}]{
                \texttt{\tiny{[GM$|$GM]}}
                \texttt{Invalid format: head and tail are not empty} \\
\texttt{head: 'Съжалявам за всички грешки. Ето перфектно форматираната команда:} \\
\\ 
\texttt{**'} \\
\texttt{tail: '**} \\
\\ 
\texttt{Това е валидна команда за преместване, която ще ми помогне да избегна наказания, докато планираме следващите стъпки. Моля потвърдете дали е приета. След това ще продължа с комуникацията за координиране на целите.'} \\
            }
        }
    }
    & & \\ \\

    \theutterance \stepcounter{utterance}  
    & & & \multicolumn{2}{p{0.3\linewidth}}{
        \cellcolor[rgb]{0.9,0.9,0.9}{
            \makecell[{{p{\linewidth}}}]{
                \texttt{\tiny{[GM$|$GM]}}
                \texttt{Maximum number of penalties exceeded} \\
            }
        }
    }
    & & \\ \\

    \theutterance \stepcounter{utterance}  
    & & & \multicolumn{2}{p{0.3\linewidth}}{
        \cellcolor[rgb]{0.9,0.9,0.9}{
            \makecell[{{p{\linewidth}}}]{
                \texttt{\tiny{[GM$|$GM]}}
                \texttt{Initial states:} \\
\texttt{```} \\
\texttt{Player 1:} \\
\texttt{ 1234567} \\
\texttt{╔═╤═════╗} \\
\texttt{║◌│◌A◌◌◌║ 1} \\
\texttt{╟─┤◌T◌◌◌║ 2} \\
\texttt{║O└─┐◌◌◌║ 3} \\
\texttt{║◌◌◌├─┐◌║ 4} \\
\texttt{║◌◌S│◌│◌║ 5} \\
\texttt{╟─┬─┴─┼─╢ 6} \\
\texttt{║P│◌DM│◌║ 7} \\
\texttt{╚═╧═══╧═╝} \\
\\ 
\\ 
\texttt{Player 2:} \\
\texttt{ 1234567} \\
\texttt{╔═╤═════╗} \\
\texttt{║P│S◌D◌◌║ 1} \\
\texttt{╟─┤◌◌◌◌◌║ 2} \\
\texttt{║◌└─┐◌A◌║ 3} \\
\texttt{║◌◌M├─┐◌║ 4} \\
\texttt{║◌T◌│◌│O║ 5} \\
\texttt{╟─┬─┴─┼─╢ 6} \\
\texttt{║◌│◌◌◌│◌║ 7} \\
\texttt{╚═╧═══╧═╝} \\
\\ 
\texttt{```} \\
            }
        }
    }
    & & \\ \\

    \theutterance \stepcounter{utterance}  
    & & & \multicolumn{2}{p{0.3\linewidth}}{
        \cellcolor[rgb]{0.9,0.9,0.9}{
            \makecell[{{p{\linewidth}}}]{
                \texttt{\tiny{[GM$|$GM]}}
                \texttt{End states:} \\
\texttt{```} \\
\texttt{Player 1:} \\
\texttt{ 1234567} \\
\texttt{╔═╤═════╗} \\
\texttt{║◌│◌A◌◌◌║ 1} \\
\texttt{╟─┤◌T◌◌◌║ 2} \\
\texttt{║O└─┐◌◌◌║ 3} \\
\texttt{║◌◌◌├─┐◌║ 4} \\
\texttt{║◌◌S│◌│◌║ 5} \\
\texttt{╟─┬─┴─┼─╢ 6} \\
\texttt{║P│◌DM│◌║ 7} \\
\texttt{╚═╧═══╧═╝} \\
\\ 
\\ 
\texttt{Player 2:} \\
\texttt{ 1234567} \\
\texttt{╔═╤═════╗} \\
\texttt{║P│S◌D◌◌║ 1} \\
\texttt{╟─┤◌◌◌◌◌║ 2} \\
\texttt{║◌└─┐◌A◌║ 3} \\
\texttt{║◌◌M├─┐◌║ 4} \\
\texttt{║◌T◌│◌│O║ 5} \\
\texttt{╟─┬─┴─┼─╢ 6} \\
\texttt{║◌│◌◌◌│◌║ 7} \\
\texttt{╚═╧═══╧═╝} \\
\\ 
\texttt{```} \\
            }
        }
    }
    & & \\ \\

    \theutterance \stepcounter{utterance}  
    & & & \multicolumn{2}{p{0.3\linewidth}}{
        \cellcolor[rgb]{0.9,0.9,0.9}{
            \makecell[{{p{\linewidth}}}]{
                \texttt{\tiny{[GM$|$GM]}}
                \texttt{* success: False} \\
\texttt{* lose: False} \\
\texttt{* aborted: True} \\
\texttt{{-}{-}{-}{-}{-}{-}{-}} \\
\texttt{* Shifts: 0.00} \\
\texttt{* Max Shifts: 12.00} \\
\texttt{* Min Shifts: 6.00} \\
\texttt{* End Distance Sum: 32.45} \\
\texttt{* Init Distance Sum: 32.45} \\
\texttt{* Expected Distance Sum: 29.33} \\
\texttt{* Penalties: 17.00} \\
\texttt{* Max Penalties: 16.00} \\
\texttt{* Rounds: 0.00} \\
\texttt{* Max Rounds: 28.00} \\
\texttt{* Object Count: 7.00} \\
            }
        }
    }
    & & \\ \\

\end{supertabular}
}

\end{document}
